\honest{Minds: An Introduction}{Rob Bensinger}{2015}

You are a mind, and very few things are. You can make plans, revise beliefs, notice ladybugs, and even reason about your own reasoning. Your mind runs on a brain, which is ``a lawful thing, a thing of pattern and routine,'' and when a mental pattern serves you well, we call that ``rationality.''

Your routines come from your ancestry. Replicating molecules produced heritable differences, some differences affected reproduction, and so evolution built organisms suited to environments like their ancestors'. You are ``a mind, carved from weaker minds,'' trying to improve yourself for your own goals and not those of ``your designer, evolution.''

Our mental categories, such as belief, decision and feeling, look nothing like our physical ones. Philosophers took this as a reason to separate minds from brains, the view Gilbert Ryle called ``the dogma of the Ghost in the Machine.'' \nb{In \textsc{The Concept of Mind} (1949) Ryle says he uses the phrase ``with deliberate abusiveness,'' for the official doctrine he traces chiefly to Descartes.} Those who reject dualism ``haven't necessarily'' found a better predictive model, so our talk of rationality and bias may still be as loose as talk of spirits. I name no one who makes this mistake. Perhaps we can learn about ourselves from processes that seem purposeful but are plainly inhuman: evolution and AI. The appearance of a ghost in the machine ``is itself the hidden work of a machine.''

I then preview the book's three sequences and its closing essay on Bayes's theorem.

The second half is about Yudkowsky, whom I call a decision theorist and mathematician. I quote Yudkowsky's first post on \textsc{Overcoming Bias} to show that the work on AI drove the interest in human rationality. I explain the ``intelligence explosion'' and ``Friendly AI,'' and say plainly that little is known about when general AI might come, and that experts disagree. In passing: ``until January 2013, MIRI was named `the Singularity Institute for Artificial Intelligence'\,'' without saying what MIRI is. \nb{It is the Machine Intelligence Research Institute, where Yudkowsky was a senior researcher. MIRI released this book, and Bensinger announced it on MIRI's blog.} I quote Russell and Norvig's textbook on the difficulty of making an AI friendly. \nb{In the textbook the passage begins as a report of Yudkowsky's own work: ``Yudkowsky (2008) goes into more detail \ldots\ He asserts \ldots''}

Why AI in a book on human rationality? It gives simple illustrations, and people, experts included, are ``extraordinarily bad'' at forecasting. \nb{No source is given. Philip Tetlock's forecasting tournaments support the claim for political forecasts.} Long-term technology forecasting is an important use of Bayesian rationality, which ``can model correct reasoning even in domains where the data is scarce or equivocal.'' \nb{With scarce data, a Bayesian answer depends mostly on the starting estimate, which consistency does not fix.}

We begin, then, with what ``our own designer'' can teach us.
